\section*{第 \ref{ch:g12:limcont} 章 --- 极限与连续}

\begin{solution}{exo:g12:limcont:1}
分子分母同除以 $x^3$：
$\dfrac{2x^3 - x + 1}{5x^3 + x^2} = \dfrac{2 - 1/x^2 + 1/x^3}{5 + 1/x}
\to \dfrac{2}{5}$.

$x^2 + 3x - 1 = x^2(1 + 3/x - 1/x^2) \to +\infty$，当 $x \to -\infty$ 时，
因 $x^2 \to +\infty$，且括号项趋于 $1$。

当 $x \to 2^+$ 时，分子趋于 $3 > 0$，分母趋于 $0^+$，
故 $\dfrac{x+1}{x-2} \to +\infty$。

当 $x > 0$ 时，$-\dfrac1x \leq \dfrac{\cos x}{x} \leq \dfrac1x$，故极限
为 $0$，由夹逼定理得出。
\end{solution}

\begin{solution}{exo:g12:limcont:2}
\emph{1.} $2x^2 - 3x + 1 = (x-1)(2x-1)$（根 $x=1$ 可由
$2 - 3 + 1 = 0$ 看出）。当 $x \neq 1$ 时，约去 $x - 1$，得 $f(x) = 2x - 1$。

\emph{2.} 因此 $\lim_{x\to1} f(x) = 1$。令 $f(1) = 1$，则 $f$ 延拓为
多项式函数 $x \mapsto 2x-1$，该函数定义在 $\R$ 上，且在 $1$ 处连续。
\end{solution}

\begin{solution}{exo:g12:limcont:3}
将 $f(x) = \dfrac{3x - 1}{x + 2} = \dfrac{3(x+2) - 7}{x+2}
= 3 - \dfrac{7}{x+2}$。当 $x \to \pm\infty$ 时，$f(x) \to 3$，直线 $y = 3$
是两端的水平渐近线。当 $x \to (-2)^\pm$ 时，
$f(x) \to \mp\infty$，直线 $x = -2$ 为铅直渐近线。曲线
为双曲线，其中位于 $y=3$ 下方的部分对应 $x > -2$，位于其上方的部分对应 $x < -2$。
\end{solution}

\begin{solution}{exo:g12:limcont:4}
乘以共轭式：
\[
\sqrt{x^2 + x} - x = \frac{x^2 + x - x^2}{\sqrt{x^2+x} + x}
= \frac{x}{x\left(\sqrt{1 + 1/x} + 1\right)}
= \frac{1}{\sqrt{1 + 1/x} + 1}
\xrightarrow[x\to+\infty]{} \frac{1}{2}.
\]
同理，
\[
\frac{\sqrt{1+x} - 1}{x} = \frac{(1+x) - 1}{x\left(\sqrt{1+x} + 1\right)}
= \frac{1}{\sqrt{1+x} + 1} \xrightarrow[x\to0]{} \frac{1}{2}.
\]
\end{solution}

\begin{solution}{exo:g12:limcont:5}
\emph{1.} 由 $\sin x \geq -1$ 得 $f(x) \geq x - 1 \to +\infty$，
由比较法即得结论。然而 $f'(x) = 1 + \cos x$ 在每个奇数
倍的 $\pi$ 处为零，并且在每个区间 $\intco{A}{+\infty}$ 上，函数
在若干严格递增的区间之间交替，它是递增的，但并非
严格递增；更重要的是，并不需要其增长性：比较法已足够。

\emph{2.} 当 $x > 1$ 时，同除以 $x$：
\[
g(x) = \frac{1 + \frac{\sin x}{x}}{1 - \frac{\sin x}{x}},
\]
且由夹逼定理得 $\frac{\sin x}{x} \to 0$，
（$\abs{\sin x} \leq 1$）。于是 $g(x) \to \frac{1+0}{1-0} = 1$。
\end{solution}

\begin{solution}{exo:g12:limcont:6}
$f(x) = x^3 - 3x + 1$ 连续，且 $f'(x) = 3x^2 - 3 = 3(x-1)(x+1)$：
$f$ 在 $\intoc{-\infty}{-1}$ 上递增，在 $\intcc{-1}{1}$ 上递减，
在 $\intco{1}{+\infty}$ 上递增，其中极大值 $f(-1) = 3 > 0$ 与
极小值 $f(1) = -1 < 0$，$\lim_{-\infty} f = -\infty$，
$\lim_{+\infty} f = +\infty$。

在三个单调区间中的每一个上，$f$ 连续、严格单调，
且 $0$ 严格介于端点处的函数值或极限之间，于是
由双射定理，每个区间恰有一个零点，共三个零点。
符号检验：$f(-2) = -1 < 0 < f(-1) = 3$，故在
$\intoo{-2}{-1}$ 内有一个零点；$f(0) = 1 > 0 > f(1) = -1$，故在 $\intoo{0}{1}$ 内有一个零点；
$f(1) = -1 < 0 < f(2) = 3$，故在 $\intoo{1}{2}$ 内有一个零点。
\end{solution}

\begin{solution}{exo:g12:limcont:7}
设 $g(x) = f(x) - x$，它在 $\intcc{0}{1}$ 上连续。因
$f(0) \in \intcc{0}{1}$，故 $g(0) = f(0) \geq 0$；因 $f(1) \in \intcc{0}{1}$，
故 $g(1) = f(1) - 1 \leq 0$。由介值定理，$g$ 在
某点 $c \in \intcc{0}{1}$ 处为零，即 $f(c) = c$。
\end{solution}

\begin{solution}{exo:g12:limcont:8}
设 $u(t)$ 与 $d(t)$ 分别为上山与下山途中的位置（以距
小径底部的距离度量），$t \in \intcc{8}{20}$ 为时刻，
两者均连续，$u(8) = 0$，$u(20) = L$（即小径的长度），
$d(8) = L$，$d(20) = 0$。函数 $g = u - d$ 连续，且
$g(8) = -L \leq 0$ 与 $g(20) = L \geq 0$，故由介值定理，$g(c) = 0$ 对某一时刻 $c$ 成立，
在时刻 $c$，两次行程的位置重合。
\end{solution}

\begin{solution}{exo:g12:limcont:9}
\emph{1.} $f'(x) = 5x^4 + 3x^2 + 1 > 0$，故 $f$ 在
$\R$ 上严格递增。结合 $\lim_{\pm\infty} f = \pm\infty$ 与连续性，
由双射定理，存在唯一的实零点 $\alpha$。由 $f(0) = -1 < 0$ 及
$f(1) = 2 > 0$，得 $\alpha \in \intoo{0}{1}$。

\emph{2.} 维持区间 $\intcc{a_n}{b_n}$，使
$f(a_n) < 0 < f(b_n)$，起始取 $\intcc{0}{1}$；每一步计算
$f$ 在中点 $m$ 处的值，并保留使符号改变
仍然成立的一半。经过 $n$ 步后，区间长度为 $2^{-n}$，故中点
与 $\alpha$ 的误差不超过 $2^{-(n+1)}$。要求
$2^{-(n+1)} \leq 10^{-6}$，即 $n + 1 \geq \frac{6}{\log_{10} 2} \approx 19.93$：
取 $n = 20$ 次迭代即可。
\end{solution}

\begin{solution}{pb:g12:limcont:1}
\textbf{1.} 主导项：$\frac{3x^2}{x^2} \to 3$；
$\frac{2x}{x^2} = \frac2x \to 0$;
$x^3 - x^2 = x^3\left(1 - \frac1x\right) \to +\infty$.

\textbf{2.} $\lim_{x \to \pm\infty} f = 3$：水平
渐近线 $y = 3$；$\lim_{x \to (-2)^\pm} f = \mp\infty$：
铅直渐近线 $x = -2$。

\textbf{3.} $\frac{(x-3)(x+3)}{x - 3} = x + 3 \to 6$。又
$\frac{\sqrt{x+1} - 1}{x} = \frac{x}{x(\sqrt{x+1} + 1)} =
\frac{1}{\sqrt{x+1} + 1} \to \frac12$.

\textbf{4.} $\abs{x \sin\frac1x} \leq \abs x \to 0$，故
极限为 $0$。又
$\abs{\frac{\sin x}{x}} \leq \frac1x \to 0$：极限为 $0$。

\textbf{5.} $\frac{\sqrt{4x^2 + x}}{x} =
\sqrt{4 + \frac1x} \to \sqrt4 = 2$（平方根函数的
连续性及复合）。

\textbf{6.} $f(x) = x^3 + x$ 连续，且
$f(1) = 2 < 5 < 10 = f(2)$：由介值定理，$f(c) = 5$ 对某个
$c \in \intoo{1}{2}$ 成立。

\textbf{7.} $f'(x) = 3x^2 + 1 > 0$：函数在
$\R$ 上严格递增且连续，极限为 $\mp\infty$：由双射
定理，每个值——特别是 $5$——
恰好被取到一次。

\textbf{8.} $p(x) = x^3\left(1 + \frac ax + \frac{b}{x^2} +
\frac{c}{x^3}\right)$：当 $x \to +\infty$ 时，$p \to +\infty$；当
$x \to -\infty$ 时，$p \to -\infty$。于是在某个很大的区间
$\intcc{-M}{M}$ 上，$p(-M) < 0 < p(M)$：由介值定理得到一个根。
同样的首项论证对任意奇次数的多项式都
原样适用——奇次项在两个无穷远处
符号恰好相反。

\textbf{9.} 二分 $n$ 次后，所得区间的长度为
$2^{-n}$；$2^{-20} \approx 9.5 \times 10^{-7} < 10^{-6}$：
需 20 步。牛顿法每步使正确数字位数翻倍：从
一位正确数字出发，约 $3$--$4$ 步即可达到六位——
求根问题中的乌龟与灵缇犬。

\textbf{10.} $g$ 在 $\intcc{-1}{1}$ 上不连续：
它在 $0$ 处发散，而区间包含该点。介值定理的
结论随假设一同失效——始终要检验
连续性\emph{在整个区间上}成立，而不只是
端点处。

\textbf{11.} $g(x) = f(x) - x$ 连续；
$g(0) = f(0) \geq 0$（值落在 $\intcc{0}{1}$ 中），且
$g(1) = f(1) - 1 \leq 0$。由介值定理，$g(c) = 0$ 对某个 $c$ 成立，
即 $f(c) = c$。

\textbf{12.} 将该国理想化为线段
$\intcc{0}{1}$：$f$ 把每个地点映到皱褶地图上其正下方所示的位置；
折叠使地图仍覆盖在该国之上，
故 $f$ 的像落在 $\intcc{0}{1}$ 中，且它连续
（纸不会撕裂）。不动点 $c$ 是
恰好位于其自身像正上方的那一点——布劳威尔
定理处理的是真实的二维地图。

\textbf{13.} $g(0) + g(\pi) = \left(T(0) - T(\pi)\right) +
\left(T(\pi) - T(2\pi)\right) = T(0) - T(2\pi) = 0$。于是
$g(0)$ 与 $g(\pi)$ 互为相反数（或都为零）：无论哪种情形，
$g$ 连续，在某点 $t^* \in \intcc{0}{\pi}$ 处取零值：
$T(t^*) = T(t^* + \pi)$——赤道上某处，两个
对径点的温度相同，在每一
时刻都如此。

\textbf{14.} 将桌子旋转四分之一圈，互换
两条对角线的角色：压在地面上的那一对桌腿与
另一对中有一条悬空的桌腿，二者互换处境，
故悬空高度——以其符号计作
“特定对角线的间隙”——在
$\theta$ 与 $\theta + \frac\pi2$ 处取相反的符号。连续的地面
使 $h$ 连续，由介值定理得到一个角
$\theta^*$ 使 $h(\theta^*) = 0$：四条腿都着地。（矩形桌子或
阶梯状地面则没有这样的好运。）

\textbf{15.} 模板：构造两个欲使其相等的量的\emph{差}
（如 $f(x) - x$；
$T(t) - T(t + \pi)$；登山者的上升位置减下降位置；
桌子的带符号间隙）；检验它同时取到两种符号的值
（常借助两端的对称性）；对
连续函数应用介值定理。重合是必然的。

\textbf{16.} 待解问题：\emph{有多少个}解
（由严格单调性与双射定理回答，
即第 7 题）以及\emph{位于何处}（由二分法
或牛顿法回答，即第 9 题）。介值定理只是纯粹的存在性判据。

\textbf{17.} $3 - \frac{7}{x + 2}$：铅直渐近线
$x = -2$，水平渐近线 $y = 3$；在 $-2$ 两侧，函数
递增（$-\frac{7}{x+2}$ 在其
定义域内递增）；曲线是参考双曲线
$-\frac7x$ 向左平移 $2$、向上平移 $3$ 所得——两支
贴近渐近线，交叉成十字。

\textbf{18.} 当 $x \neq 1$ 时，$f(x) = x + 1$，故
$\lim_{x \to 1} f = 2$：令 $f(1) = 2$ 即可恢复
连续性。该间断是一个孔洞，而非跳跃——
补上这一个缺失的点，曲线即愈合：
\emph{可去间断点}。

\textbf{19.} 在相邻两整数之间的每个区间上连续；
在每个整数 $n$ 处，它从 $n - 1$ 跳到 $n$。
取 $a = 0.9$，$b = 1.1$：$\lfloor a \rfloor = 0 < \frac12 <
1 = \lfloor b \rfloor$，但 $\lfloor x \rfloor = \frac12$ 无
解——跳跃越过了 $\frac12$，而这
正是介值定理禁止连续函数做的事。

\textbf{20.} 局部层面：每一点处函数值 $=$ 极限值——
可去孔洞（第 18 题）破坏了它，但可修补，
下取整函数的跳跃（第 19 题）则无法修补。整体层面：在一个
区间上，所有中间值都能取到——根的
存在性证明（第 6--8 题）、不动点、赤道、桌子
与登山者都出现在这里。这一口号的正式名称是
\emph{介值定理}。
\end{solution}
